\documentclass[11pt]{article}
\usepackage[margin=1in]{geometry}
\usepackage{hyperref}
\usepackage{enumitem}
\usepackage{booktabs}
\usepackage{xcolor}
\title{ProbOS --- Month 2, Week 5\\Particle Filter Core}
\author{Nisong Monyimba}
\date{\today}

\begin{document}
\maketitle

\section*{Week 5 goal}
Build \texttt{ParticleFilter} -- ProbOS's first sequential Bayesian
inference tool -- and validate it against an exact closed-form
solution (the Kalman filter on a linear-Gaussian model) before ever
touching \texttt{BatteryModel2Cell}.

\section*{Day-by-day plan}

\subsection*{Day 1 -- Algorithm design}
\begin{itemize}[leftmargin=*]
  \item Bootstrap (SIR) particle filter: predict / update / resample
  \item Log-space weights throughout, exponentiate only at the point of
    use -- same underflow-avoidance discipline as
    \texttt{Distribution.log\_pdf()} from Week 1
  \item \texttt{python/src/particle\_filter.py} skeleton:
    \texttt{ParticleFilter}, \texttt{PFResult}
\end{itemize}

\subsection*{Day 2 -- predict/update/resample implementation}
\begin{itemize}[leftmargin=*]
  \item \texttt{predict()} reuses \texttt{Model.forward\_batch()}
    unmodified -- zero changes to the Week 2 \texttt{Model} ABC
  \item \texttt{update()}: log-space reweighting via a caller-supplied
    log-likelihood function
  \item \texttt{resample()}: systematic (low-variance) resampling
  \item Effective sample size (ESS) diagnostic and resample-when-low
    trigger
\end{itemize}

\subsection*{Day 3 -- Kalman filter ground truth}
\begin{itemize}[leftmargin=*]
  \item Implement the exact 1D linear-Gaussian Kalman filter recursion
    directly (not imported from any library) as validation ground
    truth
  \item \texttt{RandomWalkModel}: a minimal linear-Gaussian test model
\end{itemize}

\subsection*{Day 4 -- Validation against Kalman ground truth}
\begin{itemize}[leftmargin=*]
  \item Posterior mean RMSE against exact Kalman filter mean
  \item Posterior std within a sane range of exact Kalman filter std
  \item Confirm Monte Carlo error shrinks from $N=100$ to $N=2000$
    (convergence, not just agreement at one $N$)
\end{itemize}

\subsection*{Day 5 -- Unit tests + hardening}
\begin{itemize}[leftmargin=*]
  \item \texttt{python/tests/test\_particle\_filter.py}: construction,
    mechanics (predict/update/resample in isolation), and the full
    Kalman ground-truth validation suite
  \item mypy strict, ruff clean
\end{itemize}

\subsection*{Day 6-7 -- Study guide + commit}
\begin{itemize}[leftmargin=*]
  \item \texttt{docs/study/study\_guide.md} entries: Naesseth et al.
    (2019), Chopin \& Papaspiliopoulos (2020), Sarkka (2013)
  \item Full test suite, CI verification
\end{itemize}

\section*{What was actually accomplished (retrospective)}

\begin{itemize}[leftmargin=*]
  \item \texttt{ParticleFilter} implemented exactly as planned
  \item 20 tests: unit-level mechanics plus the critical Kalman ground
    truth validation
  \item Genuine mathematical proof, not a smoke test: posterior mean
    and std converge to the exact closed-form Kalman filter solution,
    with Monte Carlo error provably shrinking as $N$ grows
  \item This is the same ``validate against a known analytical case
    before touching the real model'' discipline established in Week 2
    (Kim 2007) applied to a new class of algorithm (inference, not
    forward simulation)
\end{itemize}

\section*{Numbers at Week 5 close}
\begin{itemize}[leftmargin=*]
  \item Python tests: 284/284 passing (264 Month 1 + 20 new)
  \item C++ tests: 13/13 passing
  \item mypy strict: 0 errors
  \item ruff: 0 warnings
\end{itemize}

\section*{Carried into Week 6}
\texttt{ParticleFilter} is validated and ready to be applied to
\texttt{BatteryModel2Cell} with synthetic noisy observations in a
later week. Week 6 turns to a different concern: binding the Week 4
C++ kernel into the Python package via pybind11, so the two
implementations stop being parallel and disconnected.

\end{document}
